%\documentstyle[11pt,titlepag]{book}
%\documentstyle[11pt,titlepag]{article}
\documentstyle[11pt,titlepag]{report}

%\input blnkmath.tex
%\newcommand{\T}{\textstyle}
%\newfont{\mib}{cmmib10 scaled 1100}
%\newfont{\CBM}{cmbsy10 scaled 1100}


%\newcommand{\CB}[1]{\mbox{\boldmath$\cal #1$\unboldmath}}
%\newcommand{\bmi}[1]{\mbox{\boldmath$\mit #1$\unboldmath}}

%\font\luctenb=palatinob at 20pt
%\font\lucten=palatino at 17pt
%\font\lucteni=palatinoi at 17pt

\oddsidemargin=14.6mm
\evensidemargin=14.6mm
\textwidth=150mm

\nofiles

\begin{document}
\thispagestyle{empty}


\begin{flushright}
{ \LARGE \bf Seismic Wave Propagation in Complex Media:}
{\huge  summing the multiple scattering series}\\[3em]
{\large \it T.M. Nestor}\\[1.5em] 
 {\large \it Research School of Earth Sciences, Australian National 
 University, \\ Canberra, ACT { 0200}, Australia}\\[30em]
{\large \it Submitted for the degree of Doctor of Philosophy, \\
\today}
\end{flushright}



\pagenumbering{roman}


\setcounter{page}{1}

\vspace*{30mm}
\begin{center}
{\bf STATEMENT}
\end{center}


\begin{center}
\begin{minipage} {80.0mm}
Except where acknowledged, the research 
reported in this thesis is the original contribution of the author.\newline
\end{minipage}
\end{center}


\begin{center}
\begin{minipage} {80.0mm}
\vspace{20mm}
Tod Maurice Nestor
\end{minipage}
\end{center}
\newpage

\vspace*{30mm}
\begin{center}
{\bf Acknowledgments}
\end{center}
\begin{center}
\begin{minipage} {80.0mm}
The author gratefully acknowledges the financial support of the Research School of Earth Sciences, 
at the Australian National University. The support was provided under the auspices of the 
John Conrad Jaeger scholarship.

$\quad$ The author is indebted to his supervisor Prof. B.L.N. Kennett for the considerable
 effort he has expended trying to improve 
the quality of the presentation of this research, to Dr. Hiroshi Takenaka for making available his 
pseudo-spectral method code, as well as his personal generosity, and Mr. D. Brown for helpful 
advice concerning porting code to the 
vp2200 supercomputer at the supercomputing facility of the A.N.U.

$\quad$ Finally, the author wishes to express a heart felt gratitude to his wife Ngoh Ngoh and 
his parents Jackie and Ray for their support throughout his candidature.
\end{minipage}
\end{center}

\begin{samepage}
\begin{abstract}
\setcounter{page}{4}


In this thesis, we extend the plane wave reflection-transmission reflectivity method of 
{\em Kennett\/}~(1974), to plane a stratified medium with superimposed lateral heterogeneity.
We have assumed that the heterogeneity 
varies in 2-D only, although the wave field may spread out in 3-D. 

 The research 
has been concentrated on $3$ different aspects of the problem, namely:
{\bf the mathematical model}. The heterogeneity is discretized as a space filling {\em tessellation\/} of 
small squares, superimposed on an otherwise plane stratified reference medium. 
The scattering induced by the heterogeneity is expressed as a 
{\em multiple scattering series\/}. To model each successive term in this series, 
we are effectively calculating a monochromatic, multi-source/multi-receiver seismogram.
The cost of these calculations is minimized by employing a set of reference medium
reflection/transmission  recursions that generalize the standard recursions
(Kennett, 1983).

The unknown wavefield is governed by a strongly singular
integro-differential operator equation, where the integro-differential kernel is the 
product of the reference medium 
Green's tensor that propagates the field along its scattering path in the reference model, 
and the differential 
coupling matrix that expresses the scattering action of the superimposed heterogeneity at each 
spatial point of the reference model.  We show how the integro-differential kernel may be 
transformed into a nonsingular, compact integral kernel, by analytically evaluating the
 nearfield singularities of the reference medium Green's tensor and approximating the 
differential coupling operator by a multiplicative  coupling operator. We refer to this 
later transformation as the pseudo-spectral approximation. In the pseudo-spectral coupling 
approximation, each small square of the lattice is modelled as a Rayleigh point scatterer.
While the standard pseudo-spectral method samples the {\em wavefield\/} at the Nyquist 
sampling rate appropriate to the dominant frequency of the source, we employ a higher sampling rate
so that the {\em medium\/} is sampled on a subwavelength scale at each frequency. 
   

\noindent{\bf the computational representation}. The transformed operator equation is compact, 
and consequently it is 
amenable to numerical discretization. The numerical discretization of the problem is 
expressed as a matrix equation for the unknown displacement and traction components of 
the wavefield on a regular lattice. 

\noindent{\bf the computational solution}. The principal research challenge of this thesis  
concerns the construction  of practical methods for the solution of the matrix equations. 
Due to the enormous dimension of these equations, direct solution methods based 
matrix factorization are infeasible. Instead, we propose to construct the wavefield by 
summing its multiple scattering series. In the thesis, we develop alternative summation schemes 
based on the {\em Generalized minimal residual method\/} of {\em Saad \& 
Schultz\/}~(1986) and a generalized version of the {\em bi-conjugate 
gradient method\/} of {\em Lanczos\/}~(1952)  that extends the method of 
{\em Lapwood \& Hudson\/}~(1975) to arbitrarily high orders of perturbation theory,
effectively performing an {\em analytic continuation\/}
in those problems where the multiple scattering series diverges.
Both of these summation techniques  
have proved very effective for the construction of the required wavefield, 
while using orders of magnitude less computational resources 
than direct matrix inversion. 

However, the method suffers from the same limitations as other {\em numerical\/} 
wavefield modeling schemes at high temporal frequencies, because the 
effective sampling rate of the numerical representation is proportional 
to the frequency of the wavefield to be modeled. However, unlike the  finite 
difference and pseudo-spectral schemes, we incorporate the source field, 
the free surface and radiation conditions analytically. Also, the spectrum of the 
coefficient matrix in our representation has a form which makes it amenable to solution 
via the iterative solution methods introduced in this thesis. While fourth order finite 
differences in the frequency domain require the solution of equations of similar size to 
the ones required in our method, the finite difference equations are more 
difficult to solve using iterative methods, due to their unfavourable spectral properties.



\end{abstract}
\end{samepage}


\tableofcontents

%\listoffigures
\end{document}

\bye
